\begin{table*}[t]
\centering
\scriptsize
\setlength{\tabcolsep}{5pt}
\begin{tabular}{@{}lrrrr@{}}
\toprule
\textbf{Store} & \textbf{$G{=}1$} & \textbf{$G{=}32$} &
\textbf{$G{=}128$} & \textbf{$G{=}512$} \\
\midrule
32k, CPU-pinned & 8.9 & 9.2 & 10.9 & 94.0 \\
32k, GPU-resident & 8.4 & 7.7 & 5.5 & $\infty$ \\
128k, CPU-pinned & 27.6 & 29.7 & 37.4 & 520.1 \\
128k, GPU-resident & 25.8 & 26.8 & 27.2 & 164.2 \\
1M, CPU-pinned & 198.2 & 266.7 & 421.3 & 302.6 \\
1M, GPU-resident & 180.2 & 183.0 & 190.6 & 574.7 \\
\bottomrule
\end{tabular}
\caption{\textbf{Measured repeated-query break-even $Q^\star$ over the full
serving grid.}
$G$ is generated tokens; a cell is the number of queries after which cumulative
CoMem time is lower than matched $j{=}0$ raw-text replay. Both arms share
selection, example, and pack; selection cancels in
$Q^\star=(W_{\mathrm{CoMem}}-I_{j=0})/
(T_{j=0}-T_{\mathrm{CoMem}})$. The 32k CPU cell uses
$W_{\mathrm{CoMem}}{=}2.253$\,s and a 2.16\,ms raw-token index build.
Values derive from three process medians and include store fetch, model Read,
and decode. The 18 per-process records and aggregate are released with
SHA-256 manifests. $\infty$ means CoMem's measured per-query time is not
lower, so no finite query count amortizes its one-time Write in that cell.
Large-$G$ values are sensitive because decode leaves a small per-query
denominator; they are workload-specific measurements rather than a monotonic
scaling law.}
\label{tab:serving-crossover}
\end{table*}
